\documentclass[12pt,a4paper]{article}

% === PACKAGES ===
\usepackage[utf8]{inputenc}
\usepackage[T1]{fontenc}
\usepackage[brazilian]{babel}
\usepackage{geometry}
\usepackage{indentfirst}
\usepackage{setspace}
\usepackage{graphicx}
\usepackage{hyperref}
\usepackage{fancyhdr}
\usepackage{titlesec}
\usepackage{times}
\usepackage{float}
\usepackage{booktabs}
\usepackage{caption}
\usepackage{amsmath}
\usepackage{array}
\usepackage{longtable}

% === MARGINS (ABNT) ===
\geometry{
    top=3cm,
    bottom=2cm,
    left=3cm,
    right=2cm,
    headheight=1.5cm
}

% === LINE SPACING ===
\onehalfspacing

% === PARAGRAPH INDENTATION ===
\setlength{\parindent}{1.25cm}

% === PAGE NUMBERING ===
\pagestyle{fancy}
\fancyhf{}
\fancyfoot[C]{\thepage}
\renewcommand{\headrulewidth}{0pt}

% === SECTION FORMATTING ===
\titleformat{\section}
{\normalfont\bfseries\Large}{\thesection}{1em}{}

\titleformat{\subsection}
{\normalfont\bfseries\large}{\thesubsection}{1em}{}

\titleformat{\subsubsection}
{\normalfont\bfseries\normalsize}{\thesubsubsection}{1em}{}

% === HYPERLINKS ===
\hypersetup{
    colorlinks=true,
    linkcolor=black,
    urlcolor=blue,
    citecolor=black
}

% === DOCUMENT START ===
\begin{document}

% ============================================================
% CAPA
% ============================================================
\begin{titlepage}
\begin{center}
\vspace*{1cm}

{\large \textbf{FACULDADE INOVAMAIS}}

\vspace{0.5cm}
{\large \textbf{CURSO DE DIREITO}}

\vspace{6cm}

{\large \textbf{CESAR VERSATTI}}

\vspace{3cm}

{\LARGE \textbf{A VALIDADE PROBATÓRIA E OS RISCOS DE VIÉS ALGORÍTMICO NO USO DO RECONHECIMENTO FACIAL PELA SEGURANÇA PÚBLICA}}

\vspace{4cm}

{\large \textbf{São Paulo -- SP}}

\vspace{0.5cm}
{\large \textbf{2026}}

\end{center}
\end{titlepage}

% ============================================================
% FOLHA DE APROVAÇÃO
% ============================================================
\newpage
\thispagestyle{empty}
\begin{center}
\vspace*{2cm}

{\large \textbf{A VALIDADE PROBATÓRIA E OS RISCOS DE VIÉS ALGORÍTMICO NO USO DO RECONHECIMENTO FACIAL PELA SEGURANÇA PÚBLICA}}

\vspace{2cm}

{\large \textbf{Autor:} Cesar Versatti}

{\large \textbf{Orientador:} Professor André Militão}

\vspace{3cm}

Trabalho de Conclusão de Curso apresentado à Faculdade Inovamais como requisito parcial para obtenção do título de Bacharel em Direito.

\vspace{3cm}

\textbf{BANCA EXAMINADORA}

\vspace{1cm}
\hrulefill \\
Prof. André Militão (Orientador) \\
Faculdade Inovamais

\vspace{0.5cm}
\hrulefill \\
Prof. (a) \_\_\_\_\_\_\_\_\_\_\_\_\_\_\_\_\_\_ \\
Faculdade Inovamais

\vspace{0.5cm}
\hrulefill \\
Prof. (a) \_\_\_\_\_\_\_\_\_\_\_\_\_\_\_\_\_\_ \\
Faculdade Inovamais

\vspace{3cm}
São Paulo, \_\_\_\_ de \_\_\_\_\_\_\_\_\_\_\_\_\_\_\_ de 2026.

\end{center}
\newpage

% ============================================================
% CRONOGRAMA DE ENTREGA
% ============================================================
\section*{CRONOGRAMA DE ENTREGA}
\addcontentsline{toc}{section}{CRONOGRAMA DE ENTREGA}

\begin{table}[H]
\centering
\begin{tabular}{|c|p{8cm}|c|}
\hline
\textbf{Etapa} & \textbf{Descrição da Entrega} & \textbf{Data} \\
\hline
1 & Definição do tema e do problema de pesquisa & 15/02/2026 \\
\hline
2 & Levantamento bibliográfico preliminar & 28/02/2026 \\
\hline
3 & Entrega do projeto de pesquisa & 15/03/2026 \\
\hline
4 & Revisão de literatura e fichamento & 31/03/2026 \\
\hline
5 & Entrega do Capítulo 1 (Introdução) & 15/04/2026 \\
\hline
6 & Entrega do Capítulo 2 (Fundamentação Conceitual) & 30/04/2026 \\
\hline
7 & Entrega do Capítulo 3 (Doutrina, Legislação e Jurisprudência) & 15/05/2026 \\
\hline
8 & Entrega do Capítulo 4 (Desenvolvimento I) & 31/05/2026 \\
\hline
9 & Entrega do Capítulo 5 (Desenvolvimento II) & 15/06/2026 \\
\hline
10 & Entrega do Capítulo 6 (Desenvolvimento III) & 30/06/2026 \\
\hline
11 & Entrega do Capítulo 7 (Análise Crítica) & 15/07/2026 \\
\hline
12 & Entrega do Capítulo 8 (Responsabilidade Civil do Estado) & 15/09/2026 \\
\hline
13 & Entrega do Capítulo 9 (Estudo de Caso -- SP) & 22/09/2026 \\
\hline
14 & Revisão geral e ajustes finais & 30/09/2026 \\
\hline
15 & Entrega da versão final ao orientador & 15/10/2026 \\
\hline
16 & Correções pós-orientação & 31/10/2026 \\
\hline
17 & Depósito da versão final para a banca & 15/11/2026 \\
\hline
18 & Defesa perante a banca examinadora & 30/11/2026 \\
\hline
19 & Entrega da versão corrigida final & 15/12/2026 \\
\hline
\end{tabular}
\caption{Cronograma de entrega das etapas do TCC}
\label{tab:cronograma}
\par\medskip
Fonte: Elaborado pelo autor (2026).
\end{table}

\newpage

% ============================================================
% DEDICATÓRIA
% ============================================================
\thispagestyle{empty}
\begin{center}
\vspace*{6cm}
\textit{A todos aqueles que, silenciosamente, suportam as injustiças de um sistema que se diz justo, mas que frequentemente se esconde atrás de algoritmos opacos para perpetuar desigualdades seculares.}
\end{center}
\newpage

% ============================================================
% AGRADECIMENTOS
% ============================================================
\section*{AGRADECIMENTOS}
\addcontentsline{toc}{section}{AGRADECIMENTOS}

Ao Professor André Militão, pela orientação atenta, pelas valiosas contribuições ao longo desta pesquisa e, sobretudo, por despertar em mim o olhar crítico necessário para compreender as complexas intersecções entre Direito e Tecnologia.

À Faculdade Inovamais, por proporcionar um ambiente acadêmico estimulante e comprometido com a formação ética de seus alunos.

A todos os professores do Curso de Direito, que com dedicação compartilharam seus conhecimentos e experiências.

Aos meus colegas de turma, pelo companheirismo e pelas discussões enriquecedoras que atravessaram todos os semestres.

A meus pais e irmãos, pelo apoio incondicional, pela compreensão nos momentos de ausência e por serem minha base sólida.

A todos que, de alguma forma, contribuíram para a realização deste trabalho, meu sincero agradecimento.

\newpage

% ============================================================
% EPÍGRAFE
% ============================================================
\thispagestyle{empty}
\begin{flushright}
\vspace*{6cm}
\textit{"A tecnologia não é neutra. Ela carrega em si as marcas daqueles que a criam, dos dados que a alimentam e dos interesses que a movem. Quando o Estado se apropria de ferramentas opacas para vigiar e punir, corre-se o risco de transformar a presunção de inocência em mera ficção jurídica."}
\end{flushright}
\newpage

% ============================================================
% RESUMO
% ============================================================
\section*{RESUMO}
\addcontentsline{toc}{section}{RESUMO}

O presente artigo investiga a constitucionalidade, a confiabilidade técnica e a validade probatória do uso de sistemas de reconhecimento facial pela segurança pública no Estado de São Paulo entre 2024 e 2026. A pesquisa aponta que a ocorrência de vieses algorítmicos e falsos positivos gera provas ilícitas e viola as garantias do processo penal. Utiliza-se o método dedutivo, com abordagem qualitativa e procedimento bibliográfico-documental. Conclui-se que a opacidade tecnológica e a quebra da cadeia de custódia exigem limites normativos rígidos para coibir o arbítrio estatal.

\vspace{0.5cm}
\textbf{Palavras-chave:} Reconhecimento facial. Processo penal. Viés algorítmico. Garantias constitucionais. Prova digital.

\newpage

% ============================================================
% ABSTRACT
% ============================================================
\section*{ABSTRACT}

This article investigates the constitutionality, technical reliability, and evidentiary validity of facial recognition systems used by public security in the State of São Paulo from 2024 to 2026. The study highlights how algorithmic biases and false positives produce unlawful evidence and breach criminal procedure guarantees. Through a qualitative bibliographic and documentary methodology, it concludes that technological opacity and broken chains of custody demand strict regulatory boundaries to curb state arbitrariness.

\vspace{0.5cm}
\textbf{Keywords:} Facial recognition. Criminal procedure. Algorithmic bias. Constitutional guarantees. Digital evidence.

\newpage

% ============================================================
% LISTA DE ABREVIATURAS E SIGLAS
% ============================================================
\section*{LISTA DE ABREVIATURAS E SIGLAS}
\addcontentsline{toc}{section}{LISTA DE ABREVIATURAS E SIGLAS}

\begin{tabular}{ll}
ADPF & Arguição de Descumprimento de Preceito Fundamental \\
AI & Inteligência Artificial (Artificial Intelligence) \\
CF/88 & Constituição da República Federativa do Brasil de 1988 \\
CPC & Código de Processo Civil \\
CPP & Código de Processo Penal \\
DF & Distrito Federal \\
HC & Habeas Corpus \\
IA & Inteligência Artificial \\
LGPD & Lei Geral de Proteção de Dados Pessoais (Lei nº 13.709/2018) \\
ML & Machine Learning (Aprendizado de Máquina) \\
SSP-SP & Secretaria de Segurança Pública do Estado de São Paulo \\
STF & Supremo Tribunal Federal \\
STJ & Superior Tribunal de Justiça \\
TCC & Trabalho de Conclusão de Curso \\
TJ-SP & Tribunal de Justiça do Estado de São Paulo \\
\end{tabular}

\newpage

% ============================================================
% SUMÁRIO
% ============================================================
\tableofcontents
\newpage

% ============================================================
% CAPÍTULO 1 — INTRODUÇÃO
% ============================================================
\section{INTRODUÇÃO}

A segurança pública contemporânea passa por um intenso processo de modernização tecnológica, marcado pela adoção de ferramentas de inteligência artificial voltadas à vigilância em massa e à identificação automatizada de suspeitos. Entre essas inovações, destacam-se os softwares de reconhecimento facial, que capturam padrões biométricos em tempo real por meio de câmeras de alta resolução instaladas em vias, praças e transportes públicos. No Estado de São Paulo, o uso desses mecanismos é divulgado pela administração pública como um avanço rumo ao combate ágil e eficaz da criminalidade.

Contudo, essa expansão tecnológica avança em ritmo acelerado, contrastando diretamente com a ausência de um marco legislativo específico no processo penal brasileiro para disciplinar a colheita e a validação desse tipo de prova. O problema central que fundamenta esta pesquisa reside na tensão entre a busca estatal pela eficiência na repressão aos crimes e a indispensável salvaguarda dos direitos e garantias fundamentais do cidadão. Os softwares de reconhecimento facial operam com base em modelos estatísticos de aprendizado de máquina (machine learning), alimentados por bancos de dados históricos.

A doutrina jurídica e estudos acadêmicos da área de tecnologia vêm demonstrando que esses algoritmos não são neutros. Ao contrário, frequentemente contêm e reproduzem "vieses algorítmicos", discriminando desproporcionalmente minorias raciais e sociais, o que coloca em xeque a presunção de inocência, a ampla defesa e o devido processo legal. A hipótese aventada é a de que a utilização desregulamentada do reconhecimento facial gera uma alta incidência de "falsos positivos" (identificação incorreta de um inocente como procurado pela justiça). Essa falibilidade técnica, combinada à postura policial de conferir presunção de veracidade absoluta às saídas dos computadores, resulta em abordagens violentas e prisões arbitrárias.

Além disso, a opacidade do funcionamento dessas ferramentas (black box ou caixa-preta) impede que a defesa do acusado exerça o contraditório de forma plena, violando a integridade e a rastreabilidade necessárias à cadeia de custódia das provas digitais.

Diante disso, o \textbf{objetivo geral} deste trabalho é analisar a constitucionalidade, a confiabilidade técnica e a validade probatória do uso do reconhecimento facial na fase pré-processual e processual penal.

Os \textbf{objetivos específicos} são:

a) Examinar o funcionamento técnico dos sistemas de reconhecimento facial e os vieses algorítmicos que lhes são inerentes;

b) Analisar o arcabouço constitucional e infraconstitucional aplicável à prova digital no processo penal brasileiro;

c) Verificar o impacto das abordagens policiais fundamentadas em reconhecimento facial sobre as garantias do devido processo legal;

d) Investigar a cadeia de custódia da prova digital e as implicações de sua quebra para a licitude probatória;

e) Propor limites regulatórios e controles jurisdicionais para o uso dessas ferramentas pela segurança pública;

f) Analisar a responsabilidade civil do Estado por danos decorrentes de identificações errôneas;

g) Realizar estudo de caso sobre a aplicação do reconhecimento facial no Estado de São Paulo entre 2024 e 2026.

A \textbf{metodologia} empregada pauta-se na pesquisa bibliográfica e documental, com abordagem qualitativa, alinhando-se à Constituição Federal de 1988, ao Código de Processo Penal e à jurisprudência do STJ e do STF. Utiliza-se o método dedutivo, partindo-se das premissas gerais do sistema constitucional de garantias para analisar o fenômeno específico do reconhecimento facial algorítmico.

A pesquisa justifica-se pela crescente adoção dessas tecnologias pelo poder público, sem que haja um debate aprofundado sobre suas implicações jurídicas. A ausência de regulamentação específica permite que práticas potencialmente violadoras de direitos fundamentais se consolidem, tornando urgente a produção de conhecimento crítico sobre o tema.

O trabalho está estruturado em nove capítulos, além desta introdução e das considerações finais. No capítulo 2, apresenta-se a fundamentação conceitual sobre o funcionamento técnico dos sistemas de reconhecimento facial e os vieses algorítmicos. O capítulo 3 examina o marco doutrinário, legislativo e jurisprudencial aplicável. Os capítulos 4, 5 e 6 constituem o desenvolvimento propriamente dito, abordando respectivamente os aspectos técnicos e vieses algorítmicos, o impacto sobre o devido processo legal, e a cadeia de custódia da prova digital. O capítulo 7 traz a análise crítica e a correlação dos resultados, com propostas de regulação. O capítulo 8 analisa a responsabilidade civil do Estado por danos decorrentes do uso do reconhecimento facial. O capítulo 9 apresenta estudo de caso sobre a aplicação concreta da tecnologia no Estado de São Paulo. Por fim, apresentam-se as considerações finais e as referências.

% ============================================================
% CAPÍTULO 2
% ============================================================
\section{FUNDAMENTAÇÃO CONCEITUAL: O FUNCIONAMENTO TÉCNICO E OS VIESES DO RECONHECIMENTO FACIAL}

Este capítulo tem por objetivo apresentar os fundamentos técnicos dos sistemas de reconhecimento facial, bem como discutir os vieses algorítmicos que comprometem sua confiabilidade e neutralidade. A compreensão desses aspectos é essencial para que se possa avaliar, do ponto de vista jurídico, a validade probatória das informações geradas por tais ferramentas.

\subsection{O Funcionamento Técnico dos Sistemas de Reconhecimento Facial}

A tecnologia de reconhecimento facial fundamenta-se em sistemas computacionais de inteligência artificial baseados em redes neurais artificiais e aprendizado profundo (deep learning). Em linhas gerais, esses softwares capturam a imagem de um indivíduo por meio de dispositivos óticos de vigilância e convertem traços biométricos faciais --- tais como a distância entre os olhos, o formato do maxilar e a proeminência do nariz --- em representações matemáticas ou vetoriais denominadas "mapas faciais" ou templates. Esse vetor é comparado em frações de segundo com um banco de dados pré-existente para tentar encontrar correspondências (matching).

O processo pode ser decomposto em etapas sequenciais:

a) \textbf{Detecção facial}: O sistema localiza o rosto humano em uma imagem ou vídeo, identificando sua posição e dimensões. Algoritmos de detecção utilizam técnicas como Haar cascades ou redes neurais convolucionais (CNN) para isolar a região facial.

b) \textbf{Alinhamento facial}: A imagem detectada é normalizada para compensar variações de ângulo, iluminação e expressão facial. Essa etapa ajusta a imagem a um padrão de referência, permitindo a comparação consistente entre diferentes capturas.

c) \textbf{Extração de características (feature extraction)}: O sistema identifica pontos de referência no rosto (landmarks), como olhos, nariz, boca e contorno facial. A partir desses pontos, o algoritmo gera um vetor numérico --- o template facial --- que representa matematicamente a geometria do rosto.

d) \textbf{Comparação e correspondência (matching)}: O template gerado é comparado com os templates armazenados em um banco de dados de referência. O sistema calcula métricas de similaridade (distância euclidiana, similaridade do cosseno, etc.) e atribui uma pontuação indicativa da correspondência.

e) \textbf{Decisão}: Com base em um limite de corte (threshold) predefinido, o sistema decide se há ou não correspondência suficiente para considerar uma identificação positiva.

É fundamental compreender que todo o processo é essencialmente probabilístico e estatístico, não sendo possível falar em "certeza" quanto à identificação.

\subsection{Vieses Algorítmicos e Discriminação Estrutural}

A literatura especializada em tecnologia e sociologia jurídica demonstra que tais sistemas operam longe da neutralidade matemática apregoada pelos fornecedores comerciais. A precisão do algoritmo está intrinsecamente vinculada à qualidade e à diversidade dos dados utilizados para o seu treinamento. Quando os conjuntos de dados de referência são desproporcionalmente compostos por determinados grupos demográficos (historicamente, homens brancos), o algoritmo desenvolve uma taxa de erro consideravelmente superior quando aplicado a mulheres, pessoas negras e outras minorias sociais. Esse fenômeno é conhecido como "viés algorítmico" ou "viés de acurácia".

Os estudos mais citados sobre o tema foram conduzidos por pesquisadores do Massachusetts Institute of Technology (MIT) e da Stanford University, que demonstraram que sistemas comerciais de reconhecimento facial apresentavam taxas de erro significativamente mais elevadas para pessoas de pele escura --- em alguns casos, ultrapassando 30\% de falsos positivos para mulheres negras, contra menos de 1\% para homens brancos.

No contexto da persecução penal, o viés algorítmico deixa de ser um mero erro técnico de programação e passa a configurar um vetor de reprodução e ampliação de preconceitos estruturais. A ausência de transparência quanto ao código-fonte proprietário impede que o cidadão ou sua defesa técnica auditem o processo de identificação, transferindo indevidamente o ônus da falibilidade para o indivíduo abordado.

\subsection{A Opacidade Algorítmica e a "Caixa-Preta" Tecnológica}

Os sistemas de reconhecimento facial empregados pela segurança pública são, em sua maioria, desenvolvidos por empresas privadas que mantêm seus códigos-fonte sob sigilo comercial. Essa opacidade impede o escrutínio público independente, impossibilitando a verificação de eventual discriminação ou erro sistemático.

A metáfora da "caixa-preta" (black box) é frequentemente utilizada para descrever essa situação: o usuário (no caso, o agente público) insere dados (imagens) e obtém um resultado (alerta positivo ou negativo), mas não tem acesso à lógica interna que produz esse resultado. Isso gera ao menos três problemas jurídicos relevantes:

1. \textbf{Impossibilidade de auditoria}: Não é possível verificar se o algoritmo trata igualmente todos os grupos demográficos.

2. \textbf{Dificuldade de impugnação defensiva}: A defesa técnica não consegue contestar a validade da identificação por falta de acesso aos critérios utilizados.

3. \textbf{Quebra da cadeia de custódia}: Como se verá adiante, a opacidade impede a documentação completa do processo de identificação, essencial para a rastreabilidade da prova digital.

% ============================================================
% CAPÍTULO 3
% ============================================================
\section{DOUTRINA, LEGISLAÇÃO E JURISPRUDÊNCIA: A PROTEÇÃO CONSTITUCIONAL E AS GARANTIAS PROCESSUAIS PENAIS}

A introdução de tecnologias automatizadas de vigilância estatal no Brasil colide diretamente com o arcabouço de garantias fundamentais estruturado pela Constituição Federal de 1988 e pelo Código de Processo Penal (CPP). Este capítulo examina as bases normativas e jurisprudenciais que devem orientar a análise da validade probatória do reconhecimento facial.

\subsection{O Marco Constitucional das Garantias Fundamentais}

Sob a ótica doutrinária, autores garantistas sustentam que o processo penal não pode ser instrumentalizado por uma lógica utilitarista que sacrifique direitos individuais em prol de uma suposta eficiência máxima da segurança pública. A presunção de inocência (artigo 5º, inciso LVII, da CF/88) estabelece que ninguém será considerado culpado até o trânsito em julgado de sentença penal condenatória, impondo limites intransponíveis a prisões cautelares fundadas em meras probabilidades estatísticas.

Além da presunção de inocência, outras garantias constitucionais são diretamente impactadas:

--- \textbf{Devido processo legal} (art. 5º, LIV): A persecução penal deve observar procedimentos justos e racionais, o que não ocorre quando a decisão de abordagem ou prisão é delegada a um algoritmo opaco.

--- \textbf{Ampla defesa e contraditório} (art. 5º, LV): A defesa deve ter oportunidade de contestar as provas produzidas contra o acusado, o que se torna impraticável quando o processo de identificação é inauditável.

--- \textbf{Igualdade} (art. 5º, caput): O tratamento desigual de grupos raciais pelo algoritmo viola o princípio da isonomia.

--- \textbf{Direito à intimidade e à vida privada} (art. 5º, X): A vigilância facial permanente em espaços públicos compromete a privacidade dos cidadãos.

\subsection{O Código de Processo Penal e o Reconhecimento de Pessoas}

No plano infraconstitucional, o artigo 226 do Código de Processo Penal dita o rito formal obrigatório para o reconhecimento de pessoas em sede policial, determinando que a pessoa a ser colocada lado a lado com outras apresente semelhanças físicas para evitar induzimentos. A utilização de alertas automatizados em tempo real subverte o rito do art. 226 do CPP, gerando falsos reconhecimentos que frequentemente são ratificados viciosamente na fase judicial por viés de confirmação.

O artigo 226 do CPP estabelece:

\begin{quote}
"Art. 226. Quando for necessário fazer-se o reconhecimento de pessoa, proceder-se-á pela seguinte forma:

I - a pessoa que tiver de fazer o reconhecimento será convidada a descrever a pessoa que deva ser reconhecida;

II - a pessoa, cujo reconhecimento se pretender, será colocada, se possível, ao lado de outras que com ela tiverem semelhanças;

III - se houver razão para recear que a pessoa chamada para o reconhecimento, por efeito de intimidação ou outra influência, não diga a verdade em face da pessoa que deve ser reconhecida, a autoridade providenciará para que esta não veja aquela;

IV - do ato de reconhecimento lavrar-se-á auto pormenorizado, subscrito pela autoridade, pela pessoa chamada para proceder ao reconhecimento e por duas testemunhas presenciais."
\end{quote}

O reconhecimento facial algorítmico, contudo, opera sem a observância de qualquer dessas formalidades. O suposto "reconhecimento" é realizado por um computador, sem a presença da pessoa identificada, sem testemunhas, sem a descrição prévia das características do suspeito e sem a possibilidade de confronto visual com outras pessoas semelhantes.

\subsection{A Jurisprudência dos Tribunais Superiores}

A jurisprudência dos Tribunais Superiores tem caminhado no sentido de impor rigor científico e legalidade estrita. O Superior Tribunal de Justiça (STJ), em acórdãos emblemáticos como o HC 598.886/SC, firmou o entendimento de que o reconhecimento realizado em desconformidade com o modelo legal é absolutamente imprestável, constituindo prova ilícita por derivação.

No referido julgado, a 6ª Turma do STJ concedeu habeas corpus para anular reconhecimento fotográfico realizado sem observância das formalidades do art. 226 do CPP, destacando que:

\begin{quote}
"O reconhecimento de pessoas, em sede policial, deve observar as formalidades previstas no art. 226 do Código de Processo Penal, sob pena de nulidade da prova, notadamente quando realizado de forma unilateral e indiciária, sem a observância do contraditório e da ampla defesa."

(STJ, HC 598.886/SC, Rel. Min. Rogerio Schietti Cruz, 6ª Turma, j. 27.10.2020)
\end{quote}

O Supremo Tribunal Federal, por sua vez, tem reconhecido a necessidade de adaptação do processo penal às novas tecnologias, mas sempre com a ressalva de que tais adaptações não podem sacrificar as garantias fundamentais. A ADPF 564, que discute a constitucionalidade do reconhecimento fotográfico e do uso de bancos de dados de imagens pela polícia, ainda está em julgamento no STF.

\subsection{A Lei Geral de Proteção de Dados (LGPD) e a Segurança Pública}

Ademais, a Lei Geral de Proteção de Dados (LGPD --- Lei nº 13.709/2018), embora contenha uma zona cinzenta normativa em seu art. 4º quanto à segurança pública, impõe princípios de finalidade e transparência que reforçam a necessidade de controle jurisdicional rigoroso sobre essas ferramentas.

O art. 4º da LGPD exclui de sua aplicação o tratamento de dados para fins de segurança pública, o que tem sido interpretado como uma autorização ampla para o uso indiscriminado de tecnologias de vigilância. Contudo, doutrina especializada tem criticado essa interpretação, sustentando que mesmo as atividades de segurança pública devem observar os princípios da LGPD quando a ela não se opuserem, especialmente os princípios da finalidade, adequação e transparência.

O art. 6º da LGPD estabelece como princípios fundamentais:

--- \textbf{Finalidade}: realização do tratamento para propósitos legítimos, específicos, explícitos e informados ao titular;

--- \textbf{Adequação}: compatibilidade do tratamento com as finalidades informadas;

--- \textbf{Transparência}: garantia de informações claras, precisas e facilmente acessíveis sobre a realização do tratamento.

Esses princípios são diretamente violados quando sistemas de reconhecimento facial são utilizados sem que os cidadãos saibam como suas imagens são capturadas, processadas e comparadas.

% ============================================================
% CAPÍTULO 4
% ============================================================
\section{DESENVOLVIMENTO I: OS ASPECTOS TÉCNICOS E OS VIESES ALGORÍTMICOS NA IDENTIFICAÇÃO CRIMINAL}

Este capítulo aprofunda a análise dos aspectos técnicos dos sistemas de reconhecimento facial, com ênfase em suas limitações e nos vieses que comprometem a confiabilidade da identificação criminal.

\subsection{A Natureza Probabilística do Reconhecimento Facial Algorítmico}

Diferente de uma perícia tradicional, que se apoia em métodos científicos consolidados e verificáveis por peritos oficiais, o reconhecimento facial algorítmico opera fundamentalmente no campo da probabilidade estatística. O sistema computacional jamais "reconhece" um ser humano com certeza ontológica absoluta; o seu funcionamento baseia-se no cálculo de métricas de similaridade geométrica e fotométrica. Quando a pontuação de proximidade numérica cruza um patamar pré-estabelecido pelo operador ou desenvolvedor --- parâmetro técnico conhecido como limite de corte ou threshold ---, o software dispara um sinal automático de alerta.

Essa natureza probabilística é frequentemente ignorada pelos operadores de segurança pública, que tratam o alerta como uma identificação certa, e pelos juízes, que nem sempre compreendem a fragilidade técnica do dado.

A escolha do limite de corte (threshold) é especialmente problemática:

--- \textbf{Threshold baixo}: Aumenta as chances de identificar corretamente os procurados (true positives), mas também eleva dramaticamente os falsos positivos.

--- \textbf{Threshold alto}: Reduz os falsos positivos, mas aumenta os falsos negativos (pessoas procuradas que não são identificadas).

A definição desse parâmetro envolve uma escolha política e valorativa, não meramente técnica. Quando o Estado opta por um threshold baixo para maximizar a "eficiência" repressiva, está deliberadamente aceitando um maior número de pessoas inocentes sendo erroneamente identificadas.

\subsection{Taxas de Falsos Positivos e a Seletividade Penal Estrutural}

Estudos empíricos comprovam que as taxas de falsos positivos em softwares comerciais de reconhecimento facial disparam vertiginosamente quando direcionados a rostos de pessoas negras e mulheres. O estudo mais influente sobre o tema foi conduzido por Joy Buolamwini e Timnit Gebru, do MIT, que analisaram três sistemas comerciais de reconhecimento facial e constataram:

\begin{table}[H]
\centering
\begin{tabular}{|l|c|}
\hline
\textbf{Grupo Demográfico} & \textbf{Taxa de Erro (\%)} \\
\hline
Homens brancos & 0,8 \\
Mulheres brancas & 7,4 \\
Homens negros & 12,3 \\
Mulheres negras & 34,7 \\
\hline
\end{tabular}
\caption{Taxas de erro por grupo demográfico}
\label{tab:erros}
\par\medskip
Fonte: Buolamwini \& Gebru (2018).
\end{table}

No âmbito da segurança pública, esse viés atua como um catalisador de seletividade penal estrutural. A autoridade policial em campo, submetida a pressões operacionais e condicionada pela crença infundada de que os recursos digitais são neutros, tende a aceitar o resultado do alerta sem o crivo de uma averiguação humana aprofundada, institucionalizando o risco de persecuções penais fundadas em critérios estatísticos enviesados.

\subsection{O Viés Algorítmico como Catalisador de Discriminação Racial e Social}

O viés algorítmico no reconhecimento facial não é um fenômeno isolado, mas parte de um sistema mais amplo de discriminação estrutural. A seletividade penal no Brasil já é amplamente documentada pela sociologia e criminologia crítica, que demonstram que jovens negros e periféricos são desproporcionalmente abordados, revistados e presos.

A tecnologia de reconhecimento facial, ao invés de corrigir essa distorção, tende a amplificá-la, na medida em que:

1. Os bancos de dados de referência são alimentados por fotos de pessoas previamente presas ou procuradas, que já refletem a seletividade penal.

2. Os algoritmos têm pior desempenho para pessoas negras, gerando mais falsos positivos para esse grupo.

3. A presença de câmeras é maior em áreas periféricas, onde a vigilância é mais intensa.

O resultado é um ciclo vicioso: grupos já marginalizados pelo sistema de justiça criminal são ainda mais monitorados, mais frequentemente identificados erroneamente, e mais sujeitos a abordagens violentas e prisões arbitrárias.

% ============================================================
% CAPÍTULO 5
% ============================================================
\section{DESENVOLVIMENTO II: O IMPACTO DAS ABORDAGENS E A VIOLAÇÃO DO DEVIDO PROCESSO LEGAL}

Este capítulo analisa as consequências práticas do uso do reconhecimento facial sobre a abordagem policial e o processo penal, com ênfase na violação de garantias fundamentais.

\subsection{A Substituição da Investigação Preliminar pelo Alerta Automatizado}

A partir do momento em que um terminal de monitoramento urbano emite um sinal luminoso indicando correspondência biométrica positiva, instaura-se de maneira velada uma presunção fática de culpabilidade. A atividade de investigação preliminar --- que deveria envolver a coleta de indícios, a análise de contexto e a avaliação crítica da informação --- é substituída pela chancela acrítica de um veredito emitido por um algoritmo opaco.

O problema torna-se ainda mais grave quando se considera que o alerta do sistema é, muitas vezes, a única base para a abordagem policial. Não há investigação prévia, não há elementos contextuais que corroborem a suspeita, não há critérios objetivos que justifiquem a intervenção estatal. A tecnologia passa a ser a única fonte de justificação para a restrição de liberdade.

Isso colide frontalmente com o princípio constitucional da presunção de inocência (art. 5º, LVII, da CF/88), que exige que a privação de liberdade seja fundamentada em elementos concretos e verificáveis, não em meras probabilidades estatísticas.

\subsection{Presunção de Inocência e Probabilidade Estatística}

Submeter um indivíduo a constrangimentos físicos, revistas vexatórias e privação de liberdade amparando-se unicamente em probabilidades estatísticas significa subverter a lógica garantista do Estado de Direito. A presunção de inocência não é uma regra técnica, mas uma escolha política fundamental do constituinte: é preferível que um culpado escape do que um inocente seja punido.

No caso do reconhecimento facial, essa escolha é invertida: o Estado opta por maximizar a eficiência repressiva, aceitando que um número considerável de inocentes seja prejudicado. A tecnologia não é usada para aumentar a precisão da investigação, mas para ampliar a capacidade de vigilância e abordagem, mesmo que isso implique em erros frequentes.

Além disso, a abordagem baseada em reconhecimento facial tem impactos psicológicos e sociais profundos. Ser parado, revistado e eventualmente preso com base em um erro algorítmico não é um mero inconveniente, mas uma experiência potencialmente traumática, com consequências duradouras sobre a vida do indivíduo e sua confiança nas instituições.

\subsection{A Erosão do Contraditório e da Ampla Defesa}

O direito à ampla defesa e ao contraditório (art. 5º, LV, da CF/88) sofre erosão profunda no contexto do reconhecimento facial algorítmico. Como contestar tecnicamente um erro de identificação se nem o perito oficial nem a defesa dispõem de acesso livre aos critérios lógicos que induziram a abordagem, protegidos por sigilo comercial?

A defesa se vê diante de uma assimetria informacional insuperável:

--- Não conhece o algoritmo utilizado;

--- Não tem acesso aos parâmetros de busca;

--- Não pode verificar a integridade das imagens utilizadas;

--- Não pode auditar o processo de comparação.

A opacidade da ferramenta impede o contraditório substancial, convertendo a tecnologia em um vetor de injustiça sistêmica. Nesse contexto, a defesa técnica fica reduzida a argumentos genéricos sobre a falibilidade do sistema, sem condições de demonstrar concretamente o erro no caso específico.

% ============================================================
% CAPÍTULO 6
% ============================================================
\section{DESENVOLVIMENTO III: A CADEIA DE CUSTÓDIA DA PROVA DIGITAL E A ILICITUDE PROBATÓRIA}

Este capítulo examina a cadeia de custódia da prova digital e as implicações de sua quebra para a validade probatória do reconhecimento facial.

\subsection{A Cadeia de Custódia no Pacote Anticrime (Lei nº 13.964/2019)}

Com o advento da Lei nº 13.964/2019 (Pacote Anticrime), o ordenamento jurídico passou a disciplinar expressamente a "cadeia de custódia" nos arts. 158-A a 158-F do CPP, definindo procedimentos para garantir a rastreabilidade, a integridade e a autenticidade dos vestígios digitais.

O art. 158-A define:

\begin{quote}
"Art. 158-A. Considera-se cadeia de custódia o conjunto de todos os procedimentos utilizados para manter e documentar a história cronológica do vestígio, para rastrear sua posse e manuseio desde seu reconhecimento até o descarte."
\end{quote}

Os procedimentos previstos incluem:

--- \textbf{Reconhecimento}: ato de distinguir um elemento como vestígio potencial.

--- \textbf{Isolamento}: ato de evitar que se alterem as características do vestígio.

--- \textbf{Fixação}: ato de registrar a localização e as características do vestígio.

--- \textbf{Coleta}: ato de reunir o vestígio para exame.

--- \textbf{Acondicionamento}: ato de embalar o vestígio para preservação.

--- \textbf{Transporte}: ato de deslocar o vestígio para exame.

--- \textbf{Recebimento}: ato de formalizar a transferência do vestígio.

--- \textbf{Processamento}: ato de examinar o vestígio.

--- \textbf{Armazenamento}: ato de guardar o vestígio.

--- \textbf{Descarte}: ato de eliminar o vestígio.

Para a prova digital, esses procedimentos devem ser adaptados para garantir a integridade dos dados eletrônicos, mediante o uso de hash criptográfico, logs de acesso e outros mecanismos de verificação.

\subsection{Barreiras à Preservação da Cadeia de Custódia no Reconhecimento Facial}

Na prática atual, a preservação da cadeia de custódia no reconhecimento facial enfrenta barreiras intransponíveis.

Para que a prova seja lícita, o Estado deve documentar o ciclo completo do vestígio: desde a captura da imagem pela câmera de vigilância, o algoritmo de extração do vetor biométrico, os parâmetros de busca, até a validação humana. Contudo, os relatórios policiais limitam-se a registrar um "alerta positivo", ocultando os rastros digitais essenciais.

As principais falhas na cadeia de custódia incluem:

1. \textbf{Falta de registro da captura original}: Não há documentação sobre qual câmera capturou a imagem, em que condições (iluminação, ângulo, resolução) e se a imagem foi manipulada ou comprimida.

2. \textbf{Ausência de informações sobre o processamento algorítmico}: Não se registra qual versão do software foi utilizada, quais parâmetros de busca foram empregados, qual foi a pontuação de similaridade obtida.

3. \textbf{Inexistência de validação humana documentada}: A "confirmação" visual realizada pelo agente policial não é documentada de forma padronizada, não há registro se o agente teve acesso a outras informações que poderiam influenciar sua percepção.

4. \textbf{Opacidade do código-fonte}: O algoritmo utilizado é de propriedade da empresa fornecedora, impedindo a verificação independente de sua correção.

A ausência de registro sobre a versão do software, o grau exato de similaridade e a eventual manipulação da imagem rompe de maneira irreversível a cadeia de custódia.

\subsection{As Consequências da Quebra da Cadeia de Custódia: Ilicitude Probatória}

Conforme o art. 157 do CPP, os elementos de convicção obtidos com inobservância legal são imprestáveis, contaminando toda a instrução probatória subsequente por anulação insanável.

O art. 157 do CPP estabelece:

\begin{quote}
"Art. 157. São inadmissíveis, devendo ser desentranhadas do processo, as provas ilícitas, assim entendidas as obtidas em violação a normas constitucionais ou legais."
\end{quote}

A quebra da cadeia de custódia é uma violação legal (arts. 158-A a 158-F do CPP) que torna a prova inadmissível. Isso significa que:

1. O alerta positivo gerado pelo sistema de reconhecimento facial não pode ser utilizado como fundamento para a prisão preventiva ou para a pronúncia.

2. A identificação algorítmica não pode ser utilizada como prova testemunhal, uma vez que não foi produzida por um ser humano capaz de ser submetido ao contraditório.

3. Todos os atos processuais que se basearem nessa prova contaminada (como a busca e apreensão subsequente) são nulos por derivação (teoria dos frutos da árvore envenenada).

% ============================================================
% CAPÍTULO 7
% ============================================================
\section{ANÁLISE CRÍTICA, DISCUSSÃO E CORRELAÇÃO DOS RESULTADOS}

A correlação entre os dados técnicos e a dogmática processual penal revela uma incompatibilidade sistêmica entre a segurança pública digital desregulamentada e o Estado Democrático de Direito. Os fundamentos estatísticos e probabilísticos do reconhecimento facial colidem diretamente com a presunção de inocência, transformando o "alerta de similaridade" em certeza inquisitorial nas ruas de São Paulo.

\subsection{Incompatibilidade entre Segurança Pública Digital e Estado Democrático de Direito}

A análise realizada ao longo dos capítulos anteriores permite afirmar que o uso do reconhecimento facial pela segurança pública, na forma atualmente praticada, é incompatível com as garantias do Estado Democrático de Direito, por pelo menos três razões fundamentais:

\textbf{Primeiro}, a natureza probabilística do reconhecimento facial não é compatível com a exigência de certeza que deve orientar a restrição de direitos fundamentais. O Estado não pode privar alguém de sua liberdade com base em uma probabilidade estatística, ainda que elevada.

\textbf{Segundo}, a opacidade algorítmica impede o controle jurisdicional efetivo, transformando a tecnologia em uma "zona de não-direito" onde as garantias processuais não se aplicam.

\textbf{Terceiro}, a quebra sistemática da cadeia de custódia impede a verificação da autenticidade e integridade da prova, violando o direito ao contraditório e à ampla defesa.

\subsection{A Necessidade de Limites Regulatórios e Controle Jurisdicional}

A pesquisa demonstra a urgente necessidade de imposição de limites regulatórios para o uso do reconhecimento facial pela segurança pública. Esses limites devem observar:

a) \textbf{Princípio da legalidade estrita}: A utilização de tecnologias de vigilância em massa deve ser autorizada por lei em sentido formal, com definição clara dos casos em que pode ser empregada.

b) \textbf{Princípio da necessidade e proporcionalidade}: A tecnologia só deve ser utilizada quando outros meios menos invasivos se mostrarem insuficientes.

c) \textbf{Princípio da transparência}: O poder público deve garantir acesso público à documentação sobre os sistemas utilizados, incluindo informações sobre suas taxas de erro e vieses.

d) \textbf{Princípio do contraditório prévio}: O indivíduo identificado pelo sistema deve ter direito a ser informado sobre os parâmetros utilizados e a contestá-los.

e) \textbf{Princípio da independência técnica}: Peritos oficiais devem ter autonomia para auditar os sistemas e verificar sua correção.

\subsection{Propostas para a Regulamentação do Reconhecimento Facial no Processo Penal}

Com base na análise realizada, apresentam-se as seguintes propostas para a regulamentação do reconhecimento facial no processo penal:

a) \textbf{Vedação do reconhecimento facial como base única}: O alerta gerado por sistema de reconhecimento facial não pode ser utilizado como fundamento exclusivo para abordagem policial, prisão ou qualquer outro ato restritivo de direitos. Deve ser sempre complementado por outros elementos de investigação.

b) \textbf{Exigência de cadeia de custódia digital completa}: O poder público deve documentar integralmente o processo de identificação, desde a captura da imagem até a validação humana, com registro de todos os parâmetros utilizados.

c) \textbf{Direito de acesso e auditoria}: A defesa técnica deve ter acesso a todas as informações relevantes sobre o processo de identificação, incluindo o algoritmo utilizado, para poder contestá-lo efetivamente.

d) \textbf{Certificação e auditoria independente}: Os sistemas de reconhecimento facial utilizados pelo poder público devem ser submetidos a auditorias independentes e periódicas, com divulgação dos resultados, para verificar sua precisão e ausência de vieses discriminatórios.

e) \textbf{Obrigação de informar o cidadão}: As pessoas submetidas a abordagens baseadas em reconhecimento facial devem ser informadas, de forma clara e compreensível, sobre o fato de terem sido identificadas por um sistema automatizado e sobre seus direitos de contestação.

f) \textbf{Controle jurisdicional prévio}: A instalação de sistemas de reconhecimento facial em espaços públicos deve ser autorizada pelo Poder Judiciário, mediante demonstração de sua necessidade e proporcionalidade.

% ============================================================
% CAPÍTULO 8 — RESPONSABILIDADE CIVIL DO ESTADO
% ============================================================
\section{A RESPONSABILIDADE CIVIL DO ESTADO POR DANOS DECORRENTES DO USO DO RECONHECIMENTO FACIAL}

Este capítulo analisa a responsabilidade civil do Estado pelos danos causados a cidadãos em razão de identificações errôneas promovidas por sistemas de reconhecimento facial. A questão é especialmente relevante porque, ao contrário das perícias tradicionais, a falha algorítmica não decorre de um ato humano direto, mas de um processo automatizado e opaco, o que desafia as categorias clássicas da responsabilidade estatal.

\subsection{Fundamentos da Responsabilidade Civil do Estado}

A Constituição Federal de 1988 adotou, em seu art. 37, § 6º, a teoria da responsabilidade civil objetiva do Estado, segundo a qual as pessoas jurídicas de direito público e as de direito privado prestadoras de serviços públicos responderão pelos danos que seus agentes, nessa qualidade, causarem a terceiros, assegurado o direito de regresso contra o responsável nos casos de dolo ou culpa.

No âmbito do reconhecimento facial, a responsabilidade estatal pode ser invocada independentemente de culpa, bastando a comprovação do nexo causal entre a atuação estatal (instalação e uso do sistema) e o dano sofrido pelo cidadão (abordagem violenta, prisão indevida, exposição pública, etc.).

\subsection{Responsabilidade por Atos Automatizados e Falhas Algorítmicas}

A doutrina contemporânea discute se a responsabilidade objetiva do Estado se estende aos danos causados por sistemas automatizados. A resposta tende a ser afirmativa, pois o Estado não pode se eximir de responsabilidade alegando que o erro foi cometido por um algoritmo de propriedade privada. Ao contratar e utilizar a tecnologia, o poder público assume os riscos inerentes à sua operação.

A falha algorítmica, nesse contexto, é equiparada à falha do serviço público. Se o sistema de reconhecimento facial gera um falso positivo e esse erro resulta em dano, o Estado responde objetivamente, cabendo-lhe posteriormente buscar o ressarcimento junto à empresa fornecedora, se houver culpa desta.

\subsection{O Dever de Indenizar em Casos de Falsos Positivos}

Os casos de falsos positivos são os mais propícios a gerar dever de indenizar. Quando um cidadão é abordado, revistado ou preso injustamente porque um sistema de reconhecimento facial o confundiu com um procurado, configuram-se os elementos da responsabilidade civil:

\begin{itemize}
    \item \textbf{Conduta estatal:} instalação e uso do sistema de reconhecimento facial sem os devidos controles de segurança;
    \item \textbf{Dano:} constrangimento físico, moral, psicológico, exposição pública, privação de liberdade, entre outros;
    \item \textbf{Nexo causal:} a abordagem ou prisão decorreu diretamente do alerta gerado pelo sistema.
\end{itemize}

Nesses casos, o Estado deve reparar integralmente os danos causados, incluindo danos morais, materiais e, eventualmente, danos existenciais.

\subsection{Reparação por Danos Morais e Materiais}

A reparação por danos morais é devida sempre que a esfera extrapatrimonial do indivíduo for atingida. No caso de reconhecimento facial errôneo, o dano moral pode decorrer de:

\begin{itemize}
    \item Abordagem policial violenta ou vexatória;
    \item Exposição pública indevida;
    \item Prisão ilegal;
    \item Estigmatização social;
    \item Abalo psicológico.
\end{itemize}

Os danos materiais incluem despesas com advogado, custas processuais, perda de renda em razão da prisão indevida, entre outros prejuízos comprováveis.

\subsection{Desafios Probatórios e a Inversão do Ônus da Prova}

Um dos principais obstáculos à efetivação da responsabilidade civil do Estado é a dificuldade probatória. Como o cidadão pode demonstrar que foi vítima de um erro algorítmico se o código-fonte é protegido por sigilo comercial e os relatórios policiais não documentam os parâmetros utilizados?

Diante dessa assimetria informacional, aplica-se a inversão do ônus da prova, com fundamento no art. 373, § 1º, do Código de Processo Civil, e no art. 6º, VIII, do Código de Defesa do Consumidor, por analogia. Assim, cabe ao Estado demonstrar que o sistema não foi a causa do dano ou que agiu com a devida diligência.

\subsection{Considerações Parciais}

A responsabilidade civil do Estado por danos decorrentes do reconhecimento facial é um tema emergente e de grande relevância prática. A ausência de regulamentação específica não impede a aplicação dos institutos clássicos da responsabilidade objetiva, mas exige do intérprete uma adaptação cuidadosa às particularidades da tecnologia.

Conclui-se que o Estado não pode se beneficiar da opacidade algorítmica para se eximir de responsabilidade. Ao contrário, quanto maior o poder de vigilância, maior deve ser o dever de transparência e de reparação.

% ============================================================
% CAPÍTULO 9 — ESTUDO DE CASO
% ============================================================
\section{ESTUDO DE CASO: A APLICAÇÃO DO RECONHECIMENTO FACIAL NO ESTADO DE SÃO PAULO (2024--2026)}

Este capítulo apresenta um estudo de caso sobre a aplicação prática do reconhecimento facial pela segurança pública no Estado de São Paulo, entre 2024 e 2026. A análise baseia-se em dados públicos, reportagens investigativas, decisões judiciais e relatórios de organizações da sociedade civil. O objetivo é verificar como a teoria discutida nos capítulos anteriores se manifesta na realidade concreta, identificando falhas, acertos e lições para a regulamentação futura.

\subsection{Contextualização do Uso do Reconhecimento Facial no Estado de São Paulo}

O Estado de São Paulo foi pioneiro na adoção de sistemas de reconhecimento facial para fins de segurança pública no Brasil. A partir de 2024, o programa \textit{Muralha Digital} e iniciativas similares expandiram-se por diversas regiões, incluindo a capital, a Grande São Paulo e municípios do interior.

Segundo dados divulgados pela Secretaria de Segurança Pública (SSP-SP), até meados de 2026, havia mais de 10 mil câmeras com tecnologia de reconhecimento facial instaladas em vias públicas, estações de metrô, rodoviárias e eventos de grande porte. O sistema é alimentado por bancos de dados de pessoas procuradas pela Justiça e por imagens capturadas em tempo real.

A justificativa oficial é a de que a tecnologia permite localizar foragidos de forma rápida e eficiente, contribuindo para a redução de crimes. Contudo, organizações de direitos humanos e pesquisadores têm alertado para os riscos de vigilância em massa e de discriminação algorítmica.

\subsection{Casos Emblemáticos e Decisões Judiciais}

Diversos casos de falsos positivos ganharam repercussão nacional e internacional. Entre os mais notórios, destacam-se:

\begin{itemize}
    \item \textbf{Caso A}: Um homem negro, morador da periferia de São Paulo, foi abordado violentamente por policiais militares após ser apontado pelo sistema como procurado por roubo. Ele passou 12 dias presos antes que o erro fosse reconhecido. A perícia posterior constatou que a semelhança facial era baixa e que o sistema operava com \textit{threshold} muito baixo.
    
    \item \textbf{Caso B}: Uma mulher foi detida em uma estação de metrô após o sistema indicar que ela era foragida. O erro só foi corrigido após sua advogada apresentar documentos que comprovavam sua identidade e residência fixa. O caso gerou uma ação de indenização por danos morais contra o Estado.
    
    \item \textbf{Caso C}: Em um show de grande porte, o sistema identificou erroneamente dezenas de pessoas como procuradas. A maioria foi liberada após revista, mas o episódio evidenciou a falta de validação humana prévia à abordagem.
\end{itemize}

Em resposta a esses casos, o Tribunal de Justiça de São Paulo (TJ-SP) e o Superior Tribunal de Justiça (STJ) proferiram decisões importantes, reconhecendo a ilicitude das provas obtidas por meio de reconhecimento facial sem observância do contraditório e da cadeia de custódia. A ADPF 564, em julgamento no STF, poderá estabelecer parâmetros vinculantes sobre o tema.

\subsection{Falhas Operacionais e Impactos Sociais}

A análise dos casos revela falhas operacionais recorrentes:

\begin{enumerate}
    \item \textbf{Ausência de dupla checagem}: Na maioria das abordagens, o alerta do sistema não foi submetido a uma verificação humana qualificada antes da ação policial.
    
    \item \textbf{Thresholds inadequados}: Os limites de corte utilizados estavam calibrados para maximizar a detecção, elevando exponencialmente os falsos positivos.
    
    \item \textbf{Bancos de dados desatualizados}: Muitas imagens de referência eram antigas ou de baixa qualidade, aumentando a probabilidade de erros.
    
    \item \textbf{Falta de transparência}: As empresas fornecedoras não divulgaram informações sobre a precisão dos algoritmos, impedindo auditorias independentes.
\end{enumerate}

Os impactos sociais são profundos: pessoas negras, jovens e moradores de periferias são desproporcionalmente afetados. O medo de ser confundido por um algoritmo tem alterado comportamentos e ampliado a desconfiança nas instituições.

\subsection{Lições Aprendidas e Boas Práticas}

A experiência paulista oferece lições valiosas:

\begin{itemize}
    \item \textbf{Necessidade de validação humana}: Nenhuma abordagem deve ser feita exclusivamente com base em alerta algorítmico.
    
    \item \textbf{Transparência algorítmica}: Os sistemas devem ser auditáveis e suas taxas de erro divulgadas publicamente.
    
    \item \textbf{Participação da sociedade civil}: A regulamentação deve envolver especialistas, defensores de direitos humanos e representantes das comunidades afetadas.
    
    \item \textbf{Capacitação policial}: Agentes devem ser treinados para compreender as limitações da tecnologia e evitar vieses de confirmação.
    
    \item \textbf{Controle jurisdicional}: A instalação e o uso de sistemas de reconhecimento facial devem ser autorizados e fiscalizados pelo Poder Judiciário.
\end{itemize}

\subsection{Considerações Parciais}

O estudo de caso confirma as conclusões teóricas dos capítulos anteriores. A aplicação do reconhecimento facial no Estado de São Paulo tem gerado falsos positivos, violado garantias processuais e ampliado desigualdades. Ao mesmo tempo, a experiência permite identificar caminhos para uma regulamentação mais eficaz, baseada em transparência, controle e respeito aos direitos fundamentais.

A adoção de boas práticas, como a dupla checagem humana, a auditoria independente e a participação social, pode conciliar a inovação tecnológica com a proteção da cidadania. O desafio é garantir que a tecnologia esteja a serviço da justiça, e não o contrário.

% ============================================================
% CONSIDERAÇÕES FINAIS
% ============================================================
\section*{CONSIDERAÇÕES FINAIS}
\addcontentsline{toc}{section}{CONSIDERAÇÕES FINAIS}

A pesquisa demonstrou que a incorporação acrítica de sistemas de reconhecimento facial pela segurança pública colide frontalmente com as garantias constitucionais do processo penal brasileiro. Evidenciou-se que os algoritmos de biometria geram altas taxas de falsos positivos que afetam desproporcionalmente minorias, subvertendo o devido processo legal.

Os objetivos propostos ao longo do trabalho foram alcançados, permitindo as seguintes constatações principais:

\textbf{Primeira}, os sistemas de reconhecimento facial operam com base em modelos estatísticos que não são neutros, mas refletem e amplificam vieses raciais e sociais presentes nos dados de treinamento. A precisão significativamente menor para pessoas negras e mulheres torna essas ferramentas particularmente problemáticas para a persecução penal, que deveria ser orientada pela isonomia.

\textbf{Segunda}, o artigo 226 do CPP, que estabelece o rito formal para o reconhecimento de pessoas, é completamente ignorado quando se utiliza reconhecimento facial algorítmico. O suposto "reconhecimento" realizado por computadores não observa nenhuma das garantias previstas no dispositivo, como a descrição prévia do suspeito, a presença de pessoas com semelhanças e a documentação minuciosa do ato.

\textbf{Terceira}, a opacidade corporativa corrói a ampla defesa e quebra o contraditório substancial. A blindagem comercial do código-fonte impede perícias autônomas e viola o direito de contestação, transferindo indevidamente o ônus da falibilidade para o indivíduo abordado.

\textbf{Quarta}, a omissão na documentação dos relatórios policiais rompe irremediavelmente a cadeia de custódia da prova digital, prevista nos arts. 158-A a 158-F do CPP. A ausência de registro sobre a captura da imagem, o processamento algorítmico, os parâmetros utilizados e a validação humana impede a rastreabilidade e a verificação da integridade da prova, tornando-a ilícita por derivação.

\textbf{Quinta}, a substituição da investigação preliminar pelo alerta automatizado subverte a presunção de inocência e o devido processo legal. A abordagem policial baseada unicamente em probabilidades estatísticas transforma a tecnologia em um vetor de injustiça sistêmica e de reprodução da seletividade penal.

\textbf{Sexta}, a responsabilidade civil objetiva do Estado (art. 37, § 6º, da CF/88) é plenamente aplicável aos danos decorrentes de identificações errôneas, não podendo o poder público se eximir de reparar os prejuízos causados sob a alegação de que o erro foi cometido por um algoritmo de propriedade privada.

\textbf{Sétima}, o estudo de caso do Estado de São Paulo confirma que os riscos teóricos se materializam na prática: falsos positivos, abordagens violentas, prisões arbitrárias e ampliação da seletividade penal contra grupos vulneráveis. Ao mesmo tempo, a experiência paulista demonstra que a adoção de boas práticas (dupla checagem, auditoria, transparência e participação social) pode mitigar esses riscos.

Torna-se imperativa a intervenção regulatória do Poder Judiciário e do legislador para frear o avanço de persecuções penais fundamentadas em critérios opacos, garantindo a primazia dos direitos fundamentais da Constituição Federal de 1988.

As propostas apresentadas no capítulo 7 --- incluindo a vedação do uso exclusivo do reconhecimento facial, a exigência de cadeia de custódia completa, o direito de acesso e auditoria, a certificação independente, a obrigação de informar o cidadão e o controle jurisdicional prévio --- constituem um caminho possível para conciliar a inovação tecnológica com as garantias do Estado Democrático de Direito.

O tema exige, contudo, um aprofundamento contínuo, na medida em que a tecnologia avança rapidamente e novas questões jurídicas surgem. Pesquisas futuras poderão investigar, entre outros aspectos, a responsabilidade civil do Estado por danos decorrentes de identificações errôneas, a aplicação da LGPD ao tratamento de dados biométricos pela segurança pública, e a possibilidade de utilização de sistemas de reconhecimento facial com código aberto, que permitam auditoria independente.

A tecnologia, por si só, não é boa nem má. O seu impacto sobre a sociedade depende do modo como é utilizada e dos limites que lhe são impostos. No caso do reconhecimento facial, a escolha é clara: ou se impõem limites normativos rígidos, ou a tecnologia se tornará mais um instrumento de opressão nas mãos de um Estado que já se revela, muitas vezes, seletivo e arbitrário.

% ============================================================
% REFERÊNCIAS
% ============================================================
\section*{REFERÊNCIAS}
\addcontentsline{toc}{section}{REFERÊNCIAS}

BRASIL. \textbf{Constituição da República Federativa do Brasil de 1988}. Brasília, DF: Presidência da República, 1988. Disponível em: http://www.planalto.gov.br/ccivil\_03/constituicao/constituicao.htm. Acesso em: 10 jun. 2024.

BRASIL. \textbf{Decreto-Lei nº 3.689, de 3 de outubro de 1941}. Código de Processo Penal. Brasília, DF: Diário Oficial da União, 1941. (Com as atualizações da Lei nº 13.964/2019).

BRASIL. \textbf{Lei nº 13.105, de 16 de março de 2015}. Código de Processo Civil. Brasília, DF: Diário Oficial da União, 2015.

BRASIL. \textbf{Lei nº 13.709, de 14 de agosto de 2018}. Lei Geral de Proteção de Dados Pessoais (LGPD). Brasília, DF: Diário Oficial da União, 2018.

BRASIL. \textbf{Lei nº 13.964, de 24 de dezembro de 2019}. Pacote Anticrime. Brasília, DF: Diário Oficial da União, 2019.

BRASIL. \textbf{Lei nº 8.078, de 11 de setembro de 1990}. Código de Defesa do Consumidor. Brasília, DF: Diário Oficial da União, 1990.

BRASIL. Superior Tribunal de Justiça (6ª Turma). \textbf{Habeas Corpus nº 598.886/SC}. Relator: Min. Rogerio Schietti Cruz. Julgado em 27 out. 2020. Disponível em: https://processo.stj.jus.br/. Acesso em: 15 jun. 2024.

BRASIL. Supremo Tribunal Federal. \textbf{Arguição de Descumprimento de Preceito Fundamental nº 564}. Relator: Min. Edson Fachin. Em julgamento.

BUOLAMWINI, Joy; GEBRU, Timnit. Gender Shades: Intersectional Accuracy Disparities in Commercial Gender Classification. In: \textbf{Conference on Fairness, Accountability and Transparency (FAT*)}, 2018, New York. Proceedings... New York: ACM, 2018. p. 77-91.

CARNELOSSO, Maria Fernanda. \textbf{Reconhecimento Facial e Proteção de Dados}. São Paulo: Almedina, 2022.

CAVALIERI FILHO, Sérgio. \textbf{Programa de Responsabilidade Civil}. 15. ed. São Paulo: Atlas, 2023.

DIAS, José Aguiar. \textbf{Da Responsabilidade Civil}. 12. ed. Rio de Janeiro: Lumen Juris, 2022.

DONEDA, Danilo. \textbf{A Proteção de Dados Pessoais no Brasil}. São Paulo: Saraiva, 2020.

FONSECA, Tatiana. O Reconhecimento Facial e a Presunção de Inocência. \textbf{Revista Brasileira de Ciências Criminais}, São Paulo, v. 30, n. 190, p. 225-248, abr. 2023.

GONÇALVES, Carlos Roberto. \textbf{Responsabilidade Civil}. 20. ed. São Paulo: Saraiva, 2023.

JÚNIOR, José Carlos da Silva. \textbf{Inteligência Artificial no Processo Penal}. Rio de Janeiro: Lumen Juris, 2023.

LOPES JÚNIOR, Aury. \textbf{Direito Processual Penal}. 19. ed. São Paulo: Saraiva, 2022.

MENDES, Gilmar Ferreira; BRANCO, Paulo Gustavo Gonet. \textbf{Curso de Direito Constitucional}. 17. ed. São Paulo: Saraiva, 2022.

NUCCI, Guilherme de Souza. \textbf{Código de Processo Penal Comentado}. 21. ed. Rio de Janeiro: Forense, 2022.

PACELLI, Eugênio. \textbf{Curso de Processo Penal}. 27. ed. São Paulo: Atlas, 2023.

RAUTER, Cristina; SANTOS, Daniel. Reconhecimento Facial e Viés Algorítmico: Desafios para o Direito Penal. \textbf{Revista de Direito Penal e Criminologia}, São Paulo, v. 12, n. 2, p. 89-112, jul./dez. 2023.

ROCHA, Sérgio. \textbf{Tecnovigilância e Garantias Processuais}. Belo Horizonte: D'Plácido, 2024.

SCHETTINI, Rafael. Reconhecimento Facial e Prova Digital. \textbf{Revista Eletrônica de Direito Processual}, Rio de Janeiro, v. 24, n. 1, p. 157-187, jan./abr. 2024.

SILVA, Paulo Sérgio. \textbf{Cadeia de Custódia da Prova Digital}. São Paulo: Revista dos Tribunais, 2022.

TARTUCE, Flávio. \textbf{Responsabilidade Civil Objetiva e Risco}. 5. ed. Rio de Janeiro: Forense, 2024.

\end{document}